\section*{Цель}
Получение практических навыков организации процесса разработки
программного обеспечения в системе управления проектами Kaiten.

\section*{Задачи}
\begin{enumerate}
    \item Создать пространство проекта в облачной системе Kaiten.
    \item Разработать карточки проекта с артефактами.
    \item Организовать процесс разработки программного обеспечения (ПО) в облачной системе управления проектами Kaiten.
\end{enumerate}

\section{ОРГАНИЗАЦИЯ ПРОЕКТА ПО РЕАЛИЗАЦИИ WEB-ПРИЛОЖЕНИЯ}

\subsection{Описание предметной области и методологии}
Предметной областью разрабатываемого проекта является образовательный веб-справочник для школьников. Проект включает в себя веб-сайта для хранения справочной информации пр различным предметам (физика, информатика, русский язык) в приятном дизайне.

В качестве методологии разработки выбран Kanban. Данная методология позволяет визуализировать процесс работы, гибко управлять задачами с помощью списков (\textit{Запланированные}, \textit{В работе}, \textit{Готово}) и оперативно отслеживать прогресс на каждом этапе разработки.

\subsection{Организация проекта в системе Kaiten}
В облачной системе управления проектами Kaiten было создано рабочее пространство и доска проекта. На \cref{fig:general} представлена доска проекта со всеми разработанными задачами, распределенными по соответствующим колонкам.

\begin{figure}[H]
    \centering
    \includegraphics[width=1\textwidth]{figs/general.png}
    \caption{Доска проекта со всеми задачами}
    \label{fig:general}
\end{figure}

Для каждой задачи были созданы карточки с описанием и необходимыми атрибутами. На \cref{fig:design,fig:informatics,fig:physics,fig:search} продемонстрированы карточки задач. В них заданы различные атрибуты: сроки выполнения, ответственные участники, метки, чек-листы и прикрепленные файлы. В сумме на карточках закреплено более 9 различных атрибутов.

\begin{figure}[H]
    \centering
    \includegraphics[width=1\textwidth]{figs/desing.png}
    \caption{Карточка задачи по разработке дизайн-модели с атрибутами}
    \label{fig:design}
\end{figure}

\begin{figure}[H]
    \centering
    \includegraphics[width=1\textwidth]{figs/informatics.png}
    \caption{Карточка задачи по разработке справочника по информатике}
    \label{fig:informatics}
\end{figure}

\begin{figure}[H]
    \centering
    \includegraphics[width=1\textwidth]{figs/physics.png}
    \caption{Карточка задачи по разработке справочника по физике}
    \label{fig:physics}
\end{figure}

\begin{figure}[H]
    \centering
    \includegraphics[width=1\textwidth]{figs/search.png}
    \caption{Карточка задачи по реализации поиска по справочнику}
    \label{fig:search}
\end{figure}

Для контроля за ходом выполнения проекта были сгенерированы автоматические отчеты. На \cref{fig:table} представлен график выполнения задач в табличном виде, а на \cref{fig:timeline} — отчет в формате TIMELINE. В отчетах корректно отображены все задачи проекта, сроки их начала и завершения.

\begin{figure}[H]
    \centering
    \includegraphics[width=1\textwidth]{figs/table.png}
    \caption{Отчет о выполнении задач проекта в табличном виде}
    \label{fig:table}
\end{figure}

\begin{figure}[H]
    \centering
    \includegraphics[width=1\textwidth]{figs/timeline.png}
    \caption{Отчет о выполнении задач проекта в формате TIMELINE}
    \label{fig:timeline}
\end{figure}

\section*{Выводы}
В ходе выполнения практической работы были успешно решены поставленные задачи: создано пространство проекта в облачной системе Kaiten, разработаны карточки с артефактами и организован процесс разработки программного обеспечения по методологии Kanban. 

Назначение систем управления проектами (СУП) заключается в эффективном планировании, распределении ресурсов, визуализации рабочих процессов и контроле за ходом выполнения работ. В процессе выполнения задания были продемонстрированы широкие возможности СУП Kaiten по составлению отчетов: система позволяет автоматически генерировать наглядные графики в табличном и TIMELINE-форматах. Это существенно упрощает мониторинг сроков, анализ прогресса команды и выявление возможных задержек при реализации проекта.